% Part: turing-machines 
% Chapter: undecidability 
% Section: unsolvability-decision-problem

\documentclass[../../../include/open-logic-section]{subfiles}

\begin{document}

\olfileid[hi]{tur}{und}{uns} 
\olsection{निर्णय समस्या असमाधेय है}

\begin{thm}
\ollabel{thm:decision-prob}
निर्णय समस्या असमाधेय है: ऐसी कोई ट्यूरिंग मशीन~$D$ नहीं है जो
प्रथम-कोटि तर्कशास्त्र के !!a{sentence}~$!B$ को आगत के रूप में रखने वाली
पट्टी पर आरंभ किए जाने पर $D$ अंततः रुके और $!B$ के वैध होने पर और केवल तभी
~$1$, तथा अन्यथा $0$ निर्गत करे।
\end{thm}

\begin{proof}
मान लें कि निर्णय समस्या समाधेय है, अर्थात् मान लें कि कोई ट्यूरिंग
मशीन~$D$ है। तब हम विराम समस्या को निम्न प्रकार हल कर सकते हैं। हम ऐसी
ट्यूरिंग मशीन~$E$ बनाते हैं जो ट्यूरिंग मशीन~$M_e$ की संख्या~$e$ और
आगत~$w$ दिए जाने पर अनुरूप !!{sentence}~$!T(M_e, w) \lif !E(M_e, w)$
की संगणना करती है और पट्टी का सबसे बायाँ खाना पढ़ते हुए रुकती है। तब
मशीन $E \concat D$, आगत $e$ और $w$ दिए जाने पर पहले~$!T(M_e, w) \lif
!E(M_e, w)$ की संगणना करती और फिर उस आगत पर निर्णय-समस्या मशीन~$D$
चलाती। $D$, निर्गत~$1$ के साथ तभी और केवल तभी रुकती है जब
$!T(M_e, w) \lif !E(M_e, w)$ वैध हो, और अन्यथा~$0$ निर्गत करती है।
\olref[ver]{lem:halt-if-valid} और \olref[ver]{lem:valid-if-halt} के
अनुसार $!T(M_e, w) \lif !E(M_e, w)$ तभी और केवल तभी वैध है जब $M_e$
आगत~$w$ पर रुकती है। अतः $E\concat D$, आगत $e$ और $w$ दिए जाने पर
निर्गत~$1$ के साथ तभी और केवल तभी रुकती जब $M_e$ आगत~$w$ पर रुकती, और
अन्यथा निर्गत~$0$ के साथ रुकती। दूसरे शब्दों में, $E \concat D$ विराम
समस्या को हल कर देती। किंतु \olref[hal]{thm:halting-problem} से हम जानते
हैं कि ऐसी ट्यूरिंग मशीन का अस्तित्व नहीं हो सकता।
\end{proof}

\begin{cor}\ollabel{cor:undecidable-sat}%
यह अनिर्णेय है कि प्रथम-कोटि तर्कशास्त्र का कोई मनमाना !!{sentence}
संतुष्टनीय है या नहीं।
\end{cor}

\begin{proof}
  मान लें कि संतुष्टनीयता का निर्णय ट्यूरिंग मशीन~$S$ द्वारा किया जा
  सकता है। तब हम निर्णय समस्या को इस प्रकार हल कर सकते हैं: आगत के रूप
  में !!a{sentence}~$B$ दिए जाने पर $!B$ को एक खाना दाईं ओर खिसकाएँ।
  खाना~$1$ पर लौटें और प्रतीक~$\lnot$ लिखें।

  अब ट्यूरिंग मशीन~$S$ चलाएँ। वह अंततः पट्टी पर या तो निर्गत $1$ (यदि
  $\lnot !B$ संतुष्टनीय है) अथवा~$0$ (यदि $\lnot !B$ असंतुष्टनीय है) के
  साथ रुकती है। यदि खाना~$1$ पर~$\TMstroke$ है तो उसे मिटाएँ; यदि
  खाना~$1$ रिक्त है तो एक~$\TMstroke$ लिखें, फिर रुकें।

  यह ट्यूरिंग मशीन सदैव रुकती है और उसका निर्गत~$1$ तभी और केवल तभी है
  जब $\lnot !B$ असंतुष्टनीय हो, तथा अन्यथा~$0$ है। चूँकि $!B$ तभी और
  केवल तभी वैध है जब $\lnot !B$ असंतुष्टनीय हो, मशीन~$1$ तभी और केवल
  तभी निर्गत करती है जब $!B$ वैध हो, और अन्यथा~$0$; अर्थात् वह निर्णय
  समस्या को हल कर देती।
\end{proof}

\begin{explain}
अतः ऐसी कोई ट्यूरिंग मशीन नहीं है जो ``क्या $!B$ प्रथम-कोटि तर्कशास्त्र
का वैध !!{sentence} है?'' प्रश्न का सदैव सही ``हाँ'' अथवा ``नहीं'' उत्तर
दे। फिर भी ऐसी ट्यूरिंग मशीन \emph{है} जो सदैव सही ``हाँ'' उत्तर देती
है---किंतु उत्तर ``नहीं'' होने पर बस रुकती नहीं। यह प्रथम-कोटि
तर्कशास्त्र के सुदृढ़ता और पूर्णता प्रमेय तथा इस तथ्य से निकलता है कि
!!{derivation} का प्रभावी रूप से परिगणन किया जा सकता है।
\end{explain}

\begin{thm}
  \ollabel{thm:valid-ce}%
  प्रथम-कोटि !!{sentence} की वैधता अर्ध-निर्णेय है: ऐसी ट्यूरिंग
  मशीन~$E$ है जो प्रथम-कोटि तर्कशास्त्र के !!a{sentence}~$!B$ को आगत के
  रूप में रखने वाली पट्टी पर आरंभ किए जाने पर $E$ अंततः रुकती है और~$1$
  तभी और केवल तभी निर्गत करती है जब $!B$ वैध हो, किंतु अन्यथा नहीं रुकती।
\end{thm}

\begin{proof}
  प्रथम-कोटि तर्कशास्त्र की सभी संभावित !!{derivation} को एक प्रभावी
  कलनविधि द्वारा एक के बाद एक उत्पन्न किया जा सकता है। मशीन~$E$ ऐसा
  करती है और जब उसे यह दिखाने वाली !!a{derivation} मिलती है कि $\Proves
  !B$, तब वह निर्गत~$1$ के साथ रुकती है। सुदृढ़ता प्रमेय के अनुसार यदि
  $E$ निर्गत~$1$ के साथ रुकती है तो इसका कारण~$\Entails !B$ है। पूर्णता
  प्रमेय के अनुसार यदि $\Entails !B$, तो यह दिखाने वाली !!a{derivation}
  है कि~$\Proves !B$। चूँकि $E$ सभी संभावित !!{derivation} को क्रमबद्ध
  ढंग से उत्पन्न करती है, वह अंततः~$\Proves !B$ दिखाने वाली किसी एक को
  पाएगी और इसलिए अंततः निर्गत~$1$ के साथ रुकेगी।
\end{proof}

\end{document}
